% SPDX-License-Identifier: Apache-2.0
\section{Conclusion}
\label{sec:e-conclusion}

A companion paper merged two hardware description languages and gave the
result a syntax shaped after Rust. This paper removed every construct Rust
already provides and asked what was left.

Seventeen of twenty constructs go. What remains is five Rust modules
of traits and types and a crate of macros, not a grammar.

Two deletions are worth more than the count suggests. The rule that a
timeless body may not count cycles needs no enforcement, because a cycle
count is spelled \code{.await} and a plain function has none: a paragraph
of specification and a compiler check become a fact about the type system.
And a pipeline, a sequence and a set of parallel processes turn out to be
one construct driven three ways, so \code{pipeline}, \code{stage} and
\code{pipe} all go, and the live-range analysis the last of them existed to
feed is what an async state machine does already.

Two results were negative and were worth more still, because neither was
visible from the language reference. Arbitrary-width arithmetic does not
work on stable Rust. A unit's parallel processes cannot take a mutable
borrow of the unit, and the repair, a register being a cell that several
processes reach, describes a register better than the version that failed.

One wall stands. A signal has no value while the program runs, so it cannot
be the condition of an \code{if}. That is what a separate grammar was
buying, and a macro front end over Rust's own syntax tree was the way
to buy it back; \code{when!}, \code{case!} and \code{select!} are
that front end, and \code{\#[lower]} reads them.

Half of what this paper builds is Chisel, in Rust, and the paper says so.
The structural half was settled in 2012, and rebuilding it buys better
error messages and a configuration that is a type. Those are worth having
and would not justify a new toolchain on their own.

The other half is the bet. An \code{async fn} that is a sequence, a
pipeline or a process according to how it is driven, with operator
latency the mapping's decision, is not something Chisel can express, and
it is what both source languages existed to provide. Nineteen probes
show that Rust accepts it, and the experiment of
Section~\ref{sec:e-next} shows that it lowers: an \code{async fn} to
a pipeline whose depth follows its operators' latencies, a unit to a
module, and a RISC-V core~\cite{vreteno} to a netlist that simulates
against the trace of its own source, and runs at 100~MHz on an
Artix-7. What the lowering does not read yet is stated there, and it
is the next thing to do.

The project's second thesis declined embedding on the grounds that a full
toolchain had become affordable. That argument has not weakened. It now
points the other way: if a toolchain is affordable, build the one that
arrives with a type system, a module system, generics and a package manager
already written, and spend the effort on the part nobody has built.
